% SPDX-License-Identifier: LPPL-1.3c
% Copyright (C) 2026 Christophe Jorssen
% Author and Current Maintainer: Christophe Jorssen <christophe.jorssen@gmail.com>
% LPPL maintenance status: maintained
% This file is part of luafemm. See LICENSE and LICENSES.md for its terms.
\section{Planar electrostatic finite elements}
\label{sec:electric-fe}

Select |femm/electrostatic={model=finite elements,...}| on the picture.
The unknown is the real scalar electric potential $V$, with
$\mathbf E=-\nabla V$ and $\mathbf D=\epsilon\mathbf E$. The domain is a finite
rectangle, invariant perpendicular to the drawing. This is a planar Poisson
formulation, not the magnetic $A_z$ equation with relabelled output, and not a
three-dimensional point-charge solve. Relative permittivities are positive
scalars in the introductory examples. Floating conductors, positive symmetric
dielectric tensors, mixed/translated periodic boundaries, capacitance matrices
and Maxwell contour forces are described in Section~\ref{sec:electric-extended}.
Nonlinear dielectrics, coupled electric circuits, exterior transformations
are outside this interface. Native electric interchange is described in
Section~\ref{sec:electric-exchange}. The optional axisymmetric formulation is documented
in Section~\ref{sec:electric-axisymmetric}; named dielectric choices are listed
in Section~\ref{sec:dielectrics}. The physical path families below are shared
by both FE geometries, with the revolved normalization and axis rules stated there.
The source-correction provider reuses dielectric, conductor and boundary paths
with additional restrictions; see Section~\ref{sec:electric-hybrid}.

\subsection{Domain, mesh and operation order}
\label{sec:electric-fe-model}
The common |length unit| and |relative permittivity| specify physical length
and the background dielectric. For this model, the picture's |xmin|, |xmax|,
|ymin| and |ymax| specify the \emph{physical computational rectangle}. The common
|volume charge density| additionally specifies the background $\rho$ in
$\mathrm{C/m^3}$ (empty means zero). A dielectric region replaces both background
properties inside its fill; its default volume density is zero.

Use |/femm/electric/xmin| and the other view bounds to crop a figure without
changing that rectangle. A saved FE snapshot retains its physical box, mesh,
fields, source data, electrode names and solver diagnostics. It may be displayed
with a different view window and root TikZ transform. Repeating its nondefault
|length unit| and background |relative permittivity| is required; other physical
or numerical overrides on restoration are rejected.

The operation order is:
\begin{enumerate}
\item Select the model and rectangle before the picture starts.
\item Declare every dielectric, electrode, charge layer and exterior boundary.
      Seed contours and ordinary annotations can be added without solving.
\item Request a mesh, field, profile or numerical result. Meshing locks physical
      declarations; evaluation also solves. Changing physical path keys later
      does not modify a solved mesh and a new physical declaration fails.
\item Reuse the owned profiles and named solution in later pictures or axes.
\end{enumerate}
No |\femmsolve| is needed. The electric solver is independent of magnetic
materials, circuits, boundary tags and disk caches. A picture with a nonempty
|name| solves at picture end even if no field was requested, so its saved result
is complete. A picture without evaluation or a saved name need not solve.

\begin{key}{/femm/electrostatic/outer boundary=\meta{choice} (initially voltage)}
|voltage| holds the rectangle at |boundary potential|. |insulated| imposes
$\mathbf D\cdot\mathbf n=0$ on its untagged edges, where $\mathbf n$ points
out of the computational domain. Neither choice is an infinite-domain condition.
Explicit boundary paths override the default on their tagged edges. At least
one prescribed potential is required; a pure Neumann problem is rejected rather
than silently receiving a gauge constraint.
\end{key}
\begin{key}{/femm/electrostatic/boundary potential=\meta{number} (initially 0)}
Default exterior Dirichlet value in volts when |outer boundary=voltage|.
Changing it shifts only this boundary, not the electrode potentials. A gauge
shift requires adding the same voltage to \emph{all} prescribed potentials.
\end{key}
\begin{key}{/femm/electrostatic/depth=\meta{positive number} (initially 1)}
Extrusion length for reported total charge and energy, in model length units.
Per-length quantities and the solved field do not depend on it. This does not
introduce finite-depth end effects. For example the default centimetre unit
with |depth=2| represents a 2\,cm reporting length.
\end{key}
\begin{key}{/femm/electrostatic/mesh size=\meta{positive number} (initially .25)}
Global target triangular mesh spacing in model units. Physical path vertices
and intersections are constrained. Begin coarse, then refine for the observables
of interest. A local physical-path |mesh size| may request a finer spacing.
\end{key}
\begin{key}{/femm/electrostatic/mesh tolerance=\meta{nonnegative number} (initially 0)}
Mesher snapping tolerance in model units. Zero selects $10^{-12}$ times the
larger domain span. It must be positive after automatic selection and smaller
than one percent of |mesh size|. It is distinct from TikZ curve flattening and
from the coordinate-to-exterior snapping below.
\end{key}
\begin{key}{/femm/electrostatic/minimum angle=\meta{number} (initially 0)}
Optional strict minimum triangle angle in degrees, from zero to less than 34.
Zero disables the target. The existing constrained mesher diagnoses impossible
acute input corners and exhausted refinement budgets; it does not silently
report an unmet quality target as success.
\end{key}
\begin{key}{/femm/electrostatic/maximum area=\meta{nonnegative number} (initially 0)}
Optional maximum triangle area in squared model units. Zero disables this
quality target. It supplements |mesh size| and may increase the element count.
\end{key}
\begin{key}{/femm/electrostatic/maximum Steiner points=\meta{integer} (initially 10000)}
Refinement budget, from 0 to 1000000. Used by optional angle/area refinement;
it does not guarantee that every requested geometry can meet those targets.
\end{key}
\begin{key}{/femm/electrostatic/solver tolerance=\meta{number} (initially 1e-10)}
Relative tolerance of the reduced real SPD linear solve, strictly between
zero and one. Sparse equilibration, bounded preconditioning and original-system
residual/backward-error checks are shared with the numerical core. This is
not a bound on finite-element or exterior-domain error.
\end{key}
\begin{key}{/femm/electrostatic/maximum iterations=\meta{integer} (initially 10000)}
Linear iteration budget, from 1 to 1000000. Failure to satisfy the linear
acceptance criteria is an error; no unconverged potential is drawn.
\end{key}
\begin{key}{/femm/electrostatic/boundary snap tolerance=\meta{nonnegative number} (initially .0001)}
For TeX-captured paths only, coordinates within this distance of a rectangle
side snap to that side. Units are model coordinates. This accommodates decimal
PGF arithmetic and rotations at finite TeX precision; it does not snap interior
interfaces or change the mesher tolerance. Zero disables it. It must be smaller
than one percent of the shorter domain span. Reduce it for very small geometry,
or select a more suitable |length unit|. The direct Lua geometry API retains
its separate floating-point roundoff checks.
\end{key}

All these options are picture-scoped physical/numerical data, set before the
picture starts. Finite numeric literals and expanded numeric macros are
accepted; scientific notation is allowed. They do not configure an analytical
symmetry model. The relevant geometric and algebraic settings are retained in
owned snapshots.

\subsection{Physical TikZ paths}
\label{sec:electric-fe-paths}
These styles capture the evaluated PGF path and keep its ordinary drawing
options. One physical role per declaration is recommended; use ordinary
|draw|, |fill|, |rounded corners|, transformations and native PGF styles.
All physical vertices must lie in the finite rectangle. Each individual closed
contour must be simple; crossings between declarations are split by the mesher.

\begin{key}{/tikz/femm/electric region=\marg{options} (default \{\})}
Declare a closed, possibly compound dielectric/source region. Options belong
to |/femm/electric region|. The PGF |even odd rule| or |nonzero rule| determines
membership, including holes; holes retain earlier regions or the background.
The last declared dielectric wins in an overlap. Electrode interiors take
precedence over dielectric regions. No source is placed inside a solid conductor.
\end{key}
\begin{key}{/femm/electric region/relative permittivity=\meta{positive number} (initially \normalfont empty)}
Dimensionless $\epsilon_r$. Empty inherits the picture's background value.
The magnetic material catalogue is not an electric permittivity database.
\end{key}
\begin{key}{/femm/electric region/volume charge density=\meta{number} (initially 0)}
Signed prescribed $\rho$ in $\mathrm{C/m^3}$. This replaces the background
volume density in the region, even when only permittivity was explicitly
changed. Charge remains prescribed; it is not a conductor equilibrium law.
\end{key}
\begin{key}{/tikz/femm/electric electrode=\marg{options}}
Prescribe a voltage on one connected open or closed path. Options belong to
|/femm/electric electrode|. An open path is a zero-thickness electrode strip
extruded perpendicular to the section. A closed path uses |interior| to choose
between a solid conductor and a voltage boundary enclosing dielectric material.
Electrodes must not overlap or touch under distinct names; merge a connected
conductor under one name. Conflicting potentials fail explicitly.
An electrode on the exterior overrides the default exterior voltage. It cannot
conflict with an explicitly tagged voltage or displacement condition there.
\end{key}
\begin{key}{/femm/electric electrode/name=\meta{name} (initially \normalfont empty)}
Name used to retrieve the electrode reaction charge. Empty assigns
|electrode 1|, |electrode 2|, and so on in declaration order. Explicit names
are preferable for measurements. Repeated names connect electrode parts with identical prescribed data;
|outer boundary| is reserved for the artificial exterior's reaction.
\end{key}
\begin{key}{/femm/electric electrode/potential=\meta{number} (initially \normalfont required)}
Prescribed real potential in volts. There is no implicit voltage inferred from
the path's colour, label or name. Zero and negative values are valid.
\end{key}
\begin{key}{/femm/electric electrode/interior=\meta{choice} (initially conductor)}
For a closed path, |conductor| fills the simple loop with constant potential
and zero field. Its triangles are excluded from dielectric assembly.
|dielectric| constrains only the loop, retaining its interior dielectric/source
regions. Use this for a grounded tube around a charged cylinder, or the outer
electrode of a coaxial capacitor. An open path has no solid interior under
either choice. A compound conductor with holes is not defined by this key.
\end{key}
\begin{key}{/tikz/femm/electric charge layer=\marg{options}}
Add a prescribed sheet charge on one open or closed internal dielectric path.
Options belong to |/femm/electric charge layer|. Path direction does not change
the sign of this scalar source. Its weak edge load produces the normal-D jump.
Overlapping source layers add their densities. A layer must not lie on the
exterior or on/inside a solid or voltage electrode; use an exterior displacement
condition or a voltage constraint as appropriate.
\end{key}
\begin{key}{/femm/electric charge layer/surface charge density=\meta{number} (initially \normalfont required)}
Signed $\sigma$ in $\mathrm{C/m^2}$. The planar line represents a surface
extruded perpendicular to the section. Its charge per length is the integral
of $\sigma$ along the line measured in metres.
\end{key}
\begin{key}{/tikz/femm/electric boundary=\marg{options} (default \{\})}
Tag one connected path on the rectangular exterior.
Use the |/femm/electric boundary| family to set its options. Every resulting segment must lie on that exterior;
interior paths fail. Tagged edges override the picture's default exterior
condition. Overlapping boundary declarations and incompatible corner potentials
are rejected. Insulating and voltage sides can meet at a corner; its voltage
then fixes the shared node.
\end{key}
\begin{key}{/femm/electric boundary/condition=\meta{choice} (initially insulated)}
|insulated| sets $\mathbf D\cdot\mathbf n=0$; |voltage| sets |potential|;
|displacement| sets |normal displacement|. The normal is outward from the
computational rectangle, independently of the captured path's direction.
\end{key}
\begin{key}{/femm/electric boundary/potential=\meta{number} (initially \normalfont empty)}
Required voltage in volts for |condition=voltage|. Ignored for the other
condition choices. Different constant-voltage sides must have compatible
values at a shared corner.
\end{key}
\begin{key}{/femm/electric boundary/normal displacement=\meta{number} (initially \normalfont empty)}
Required signed outward $g=\mathbf D\cdot\mathbf n$ in $\mathrm{C/m^2}$ for
|condition=displacement|. The weak load has sign $-g$. Positive values mean
flux leaving the dielectric domain; reversing the TikZ path does not reverse it.
\end{key}
\begin{key}{/femm/electric region/mesh size=\meta{nonnegative number} (initially 0)}
Local spacing in model units; zero inherits the global spacing. A smaller
value refines the constrained segments. For a dielectric region it also seeds
the interior more densely; an open path does not define an interior refinement
patch. Values larger than the global spacing do not coarsen the global mesh.
\end{key}
\begin{key}{/femm/electric region/curve tolerance=\meta{positive number} (initially .001)}
Chord-distance bound in model units when flattening evaluated TikZ curves and
rounded corners. The original TikZ drawing retains its curves; the physical
mesh uses their polygonal approximation. Refine this separately from mesh
spacing when geometry error matters. This key has ordinary path-local scope.
\end{key}
\begin{key}{/femm/electric electrode/mesh size=\meta{nonnegative number} (initially 0)}
Local spacing in model units; zero inherits the global spacing. A smaller
value refines the constrained segments. For a dielectric region it also seeds
the interior more densely; an open path does not define an interior refinement
patch. Values larger than the global spacing do not coarsen the global mesh.
\end{key}
\begin{key}{/femm/electric electrode/curve tolerance=\meta{positive number} (initially .001)}
Chord-distance bound in model units when flattening evaluated TikZ curves and
rounded corners. The original TikZ drawing retains its curves; the physical
mesh uses their polygonal approximation. Refine this separately from mesh
spacing when geometry error matters. This key has ordinary path-local scope.
\end{key}
\begin{key}{/femm/electric charge layer/mesh size=\meta{nonnegative number} (initially 0)}
Local spacing in model units; zero inherits the global spacing. A smaller
value refines the constrained segments. For a dielectric region it also seeds
the interior more densely; an open path does not define an interior refinement
patch. Values larger than the global spacing do not coarsen the global mesh.
\end{key}
\begin{key}{/femm/electric charge layer/curve tolerance=\meta{positive number} (initially .001)}
Chord-distance bound in model units when flattening evaluated TikZ curves and
rounded corners. The original TikZ drawing retains its curves; the physical
mesh uses their polygonal approximation. Refine this separately from mesh
spacing when geometry error matters. This key has ordinary path-local scope.
\end{key}
\begin{key}{/femm/electric boundary/mesh size=\meta{nonnegative number} (initially 0)}
Local spacing in model units; zero inherits the global spacing. A smaller
value refines the constrained segments. For a dielectric region it also seeds
the interior more densely; an open path does not define an interior refinement
patch. Values larger than the global spacing do not coarsen the global mesh.
\end{key}
\begin{key}{/femm/electric boundary/curve tolerance=\meta{positive number} (initially .001)}
Chord-distance bound in model units when flattening evaluated TikZ curves and
rounded corners. The original TikZ drawing retains its curves; the physical
mesh uses their polygonal approximation. Refine this separately from mesh
spacing when geometry error matters. This key has ordinary path-local scope.
\end{key}

\begin{codeexample}[width=0pt]
\begin{tikzpicture}[femm/electrostatic={model=finite elements,
    xmin=0,xmax=2,ymin=0,ymax=1,outer boundary=insulated,mesh size=.2}]
  \draw[thick,femm/electric electrode={name=reference,potential=0}]
    (0,0)--(0,1);
  \draw[dashed,femm/electric boundary={condition=displacement,
      normal displacement=1e-9}]
    (2,0)--(2,1);
  \pic[femm/electric={seeds={.5/.2,.5/.5,.5/.8}}] {femm electric field};
  \pic {femm equipotentials};
\end{tikzpicture}
\end{codeexample}

\subsection{Views, profiles and integral results}
\label{sec:electric-fe-views}
\begin{key}{/tikz/pics/femm electric mesh=\meta{pic}}
Generate and draw the constrained mesh of an FE electric model. Meshing locks
physical declarations but does not itself solve. It is an error to use this
pic with an analytical model.
\end{key}
\begin{key}{/tikz/femm electric mesh=\meta{style}}
Default mesh style: pale grey, .15\,pt line width. Use |.append style| or native
TikZ styles to change it without altering any numerical setting.
\end{key}
\begin{key}{/tikz/femm electric dielectric=\meta{style}}
Default geometry-view dielectric style: blue-grey outline, pale blue fill,
.5\,pt line width. The saved-geometry pic draws compound contour boundaries
without filling them, to avoid filling a hole incorrectly. The original physical
TikZ path retains the user's ordinary fill rule and drawing options.
\end{key}
\begin{key}{/tikz/femm electric charge layer=\meta{style}}
Default saved-geometry sheet style: red dashed outline, .7\,pt line width.
This drawing style (with spaces) is distinct from the physical declaration
|femm/electric charge layer| (with a slash).
\end{key}
\begin{key}{/tikz/femm electric boundary=\meta{style}}
Default saved-geometry exterior style: grey densely dotted line, .5\,pt width.
The geometry pic draws the finite rectangle, declared material contours and
physical paths. Solid electrodes reuse |femm electric conductor|; open/electrode
boundaries reuse |femm electric electrode|. Geometry display does not add sources.
\end{key}

The common |femm electric geometry| pic also accepts FE snapshots. The common
field, contour and vector pics retain their ordinary styles and view bounds.
FE equipotentials are extracted from the continuous P1 potential itself:
|levels| and |level count| select values, while |contour resolution| is relevant
only to analytical providers. Avoid constant electrode/plateau values if a
unique contour curve is intended.

Field lines walk the raw constant gradient inside each triangle. They change
direction at mesh faces without averaging a dielectric or charged interface.
They terminate at conductors, view/domain edges, a null field, a non-enterable
raw interface direction, or a tracing budget. A non-enterable direction can
occur near a discrete stagnation point or at a poorly resolved interface;
refine and inspect it rather than assuming the displayed path is complete.
|step| and the RK |tolerance| affect analytical integration only; FE
|max steps| counts element crossings, and |max length| bounds each direction
from a seed in model units. Arrows always follow $\mathbf E$.

Explicit seeds and captured seed contours work unchanged. Otherwise the FE
renderer seeds voltage edges by geometric arclength. |lines per charge| is the
total seed target for this model. |seed sign=positive| selects edges where E
enters the dielectric from an electrode, |negative| the opposite direction,
and |all| both. |auto| prefers positive edges, falling back to negative ones.
Artificial voltage boundaries may also supply seeds. This is a drawing choice,
not flux-weighted sampling. For controlled illustrations use an explicit seed
contour. Vector arrows crossing physical interfaces or a conductor are omitted.

Profiles partition the flattened measurement path at actual mesh faces. They
retain continuous V and one-sided raw E/D values, separated by NaN rows where
the gradients differ. Thus their row count may greatly exceed |samples|.
Uniform extra samples do not smooth or refine a solved field. |outside=error|
rejects portions outside the FE rectangle; |outside=nan| leaves explicit gaps.
At an exact face the point sampler selects the lowest-index incident element;
crossing profiles supply both sides. There is no recovered-field option here.

\begin{key}{/femm/electric value/electrode=\meta{name} (initially \normalfont empty)}
For |\femmelectricvalue|, select a named FE electrode's charge or voltage.
|outer boundary| selects the reaction on exterior voltage nodes not owned by
a named electrode. Empty selects global FE quantities. This option is local
to the command; it is rejected on analytical models.
\end{key}

The existing |quantity| choice gains |nodes|, |triangles|, |iterations|,
|relative residual| and |charge balance per length| for FE models.
Node/triangle counts include the constant-potential triangles of solid
conductors. The residual is the reduced linear system's relative residual;
it is not a mesh-error estimate. Units and ownership are:
\begin{center}\small
\begin{tabular}{lll}
Selection & Quantity & Meaning/units\\\hline
Global & |charge per length| & Prescribed volume plus sheet charge, C/m\\
Global & |energy per length| & Dielectric field energy, J/m\\
Global & |charge|, |energy| & Above quantities times depth, C and J\\
Electrode & |charge per length|, |charge| & Variational reaction, C/m and C\\
Electrode & |voltage| & Prescribed or solved potential, V\\
Global & |charge balance per length| & Signed balance residual, C/m\\
Global & |nodes|, |triangles|, |iterations| & Dimensionless counts\\
Global & |relative residual| & Dimensionless linear residual
\end{tabular}
\end{center}
An exterior with different voltage tags has no unique |voltage| result.
There is no general FE scalar |capacitance| quantity: a charge/voltage ratio
needs a specified source-free electrode experiment and consistent reference
potentials. Use |\femmelectricmatrix| (Section~\ref{sec:electric-extended})
for Maxwell coefficients and separate source offsets. Integral reactions converge
with the physical discretization; a small algebraic balance alone does not
certify the capacitance or a local field maximum.

\begin{codeexample}[width=0pt]
\begin{tikzpicture}[femm/electrostatic={model=finite elements,
    xmin=-2.4,xmax=2.4,ymin=-2.4,ymax=2.4,mesh size=.15,name=fe coax}]
  \draw[fill=gray!15,femm/electric electrode={name=inner,potential=10}]
    (0,0) circle[radius=1];
  \draw[thick,femm/electric electrode={name=outer,potential=0,
      interior=dielectric}]
    (0,0) circle[radius=2];
  \pic {femm electric field};
  \pic[femm/electric={levels={2,4,6,8}}] {femm equipotentials};
\end{tikzpicture}
\end{codeexample}
\begin{codeexample}[width=0pt]
$Q_{\rm inner}/L=\femmelectricvalue[model=fe coax,electrode=inner,
  quantity=charge per length]\,\mathrm{C/m}$.
\quad
$W/L=\femmelectricvalue[model=fe coax,quantity=energy per length]\,\mathrm{J/m}$.
\end{codeexample}

\subsection{Weak formulation, reactions and limits}
\label{sec:electric-fe-methods}
With lengths in metres, prescribed internal sheet charge $\sigma$, and outward
exterior displacement $g$, the weak equation for free test functions is
\[
 \int_\Omega\epsilon\nabla V\cdot\nabla w\,dA
 =\int_\Omega\rho w\,dA+\int_{\Gamma_\sigma}\sigma w\,dl
   -\int_{\Gamma_N}g w\,dl.
\]
Dirichlet values are eliminated before solving. Constant coefficients and
sources are integrated exactly on straight P1 triangles and segments.
Dividing matrix and loads by $\epsilon_0$ improves numerical scale without
changing the law. Explicit model-to-SI factors remain on volume and edge loads.
The constrained triangular mesher, predicates, spatial acceleration and audited
sparse CG solver are shared infrastructure; the electric assembly, loads,
units and observables are separate from magnetic constitutive laws.

For a triangle $T$, $\mathbf E_T=-\sum_i V_i\nabla N_i$ is constant and
$W/L=\sum_T\epsilon_T|\mathbf E_T|^2|T|/2$. Solid conductor triangles contribute
neither dielectric stiffness nor field energy. Fixed-node reactions of the
\emph{full} assembled equation give electrode charge per length; they include
the volume/surface loads in the electrode-adjacent test functions. Summing
prescribed volume and sheet charge, electrode/exterior reactions, and subtracting
prescribed outward Neumann flux gives the reported charge-balance residual.
Raw-gradient surface quadrature need not equal this discrete conservative reaction.

Validation includes an exact two-dielectric slab, a charged sheet, a manufactured
volume-charge parabola, displacement signs, planar unit/depth scaling, compound
holes, snapshots and independent mesh/exterior sensitivity. A coaxial capacitor
is compared with $C/L=2\pi\epsilon/\log(b/a)$ and $2W=QU$. Curves, material
interfaces, sharp conductor edges and finite-domain truncation remain separate
sources of error. Point-charge nodes retain their three-dimensional Coulomb
meaning and are rejected in this planar FE model. Axisymmetric electrostatics
requires a different weighted scalar formulation and is not selected here.
